\documentclass[UTF8]{article}
\usepackage[UTF8]{ctex}
\usepackage{amsmath, amssymb}
\usepackage{physics}
\usepackage{graphicx}
\usepackage{subcaption}
\usepackage{caption}
\usepackage[font=small]{caption}
\usepackage{booktabs}
\usepackage{multirow}
\usepackage{geometry}
\geometry{a4paper, left=2.5cm, right=2.5cm, top=2.5cm, bottom=2.5cm}
\usepackage{enumitem}
\usepackage{float}
\usepackage{siunitx}
\sisetup{per-mode=symbol}
\usepackage[hidelinks]{hyperref}

% 页眉页脚
\usepackage{fancyhdr}
\pagestyle{fancy}
\lhead{干涉衍射实验}
\chead{}
\rhead{物理实验B}
\lfoot{}
\cfoot{\thepage}
\rfoot{}
\renewcommand{\headrulewidth}{0.4pt}
\setlength{\headheight}{14pt}

\title{波动光学——干涉与衍射实验}
\author{} % 填写你的姓名和学号
\date{} % 填写实验日期

\begin{document}

\maketitle

% ============================================================
\section{实验目的}
% ============================================================

\begin{enumerate}
    \item 了解光学平台的使用方法，掌握光路搭建与等高共轴调节技术。
    \item 了解迈克尔逊干涉装置的结构和干涉条纹的形成原理。
    \item 观察等倾和等厚干涉条纹，利用迈克尔逊干涉仪测定空气折射率。
    \item 观察单缝衍射现象，学习利用衍射测微小量和测波长的方法。
    \item 利用狭缝测定激光笔的波长。
    \item 利用测定的激光笔，测定手机屏幕的像素大小并估计手机分辨率。
\end{enumerate}

% ============================================================
\section{实验原理}
% ============================================================

\subsection{等倾干涉}

发生等倾干涉的薄膜是均匀的，光线以一定倾角入射，上下两条反射光线经过透镜汇聚，形成干涉。由于入射角相同的光经薄膜两表面反射形成的反射光在相遇点有相同的光程差，凡入射角相同的就形成同一条纹，故这些倾斜度不同的光束经薄膜反射所形成的干涉花样是一些明暗相间的同心圆环，这种干涉称为\textbf{等倾干涉}。

如图所示，1与2'的光程差为：
\begin{equation}
    \delta = 2d\cos\theta
    \label{eq:equal-incline}
\end{equation}
即入射角相同的点的光程差相同，故称等倾干涉。在无穷远处形成干涉图样为同心圆。明条纹条件为$\delta = K\lambda$，暗条纹条件为$\delta = (2K+1)\lambda/2$。

\subsection{等厚干涉}

等厚干涉是由平行光入射到厚度变化均匀、折射率均匀的薄膜上、下表面而形成的干涉条纹。薄膜厚度相同的地方形成同条干涉条纹，故称等厚干涉。在两表面倾斜角度较小时，光（波）程差可以写为：
\begin{equation}
    \delta = 2d\cos\theta \approx 2d
    \label{eq:equal-thick}
\end{equation}
即空气隙厚度相同的点的光程差相同，故称等厚干涉。形成的干涉图样为明暗相间的条纹形状。

\subsection{迈克尔逊干涉实验原理}

迈克尔逊干涉仪的原理是一束入射光分为两束后各自对应的平面镜反射回来，这两束光从而能够发生干涉。干涉中两束光的不同光程可以通过调节干涉臂长度以及改变介质的折射率来实现，从而能够形成不同的干涉图样。

考虑平面镜M1和平面镜M2之间的介质和传播路程不一样，当光束垂直入射至M1、M2镜时，光程差：
\begin{equation}
    \delta = 2(n_1 l_1 - n_2 l_2)
    \label{eq:michelson-opd}
\end{equation}
当$\delta = K\lambda$，$K=0,1,2,3,\ldots$时，接收屏中心的总光强为极大，为亮条纹。

\subsection{空气折射率的测量}

正常状态（$t=25\,^{\circ}\mathrm{C}$，$P=1.01325\times10^5\,\mathrm{Pa}$）下，532\,nm的绿激光在空气中的折射率$n=1.00027821$，它与真空折射率之差为$n-1=2.782\times10^{-4}$。用一般方法不易测出这个折射率差，而用干涉法能很方便地测量，且准确度高。

在原来干涉路径中，加入一个长度为$L$、且可以调节气压$p$的玻璃气室。理论证明，在温度和湿度一定的条件下，当气压不太大时，气体折射率的变化量$\Delta n$与气压的变化量$\Delta p$成正比：
\begin{equation}
    \frac{\Delta n}{\Delta p} = \frac{\Delta n_0}{\Delta p_0}
    \label{eq:n-p-ratio}
\end{equation}

当支路L1上介质折射率改变$\Delta n$时，引起的干涉条纹的变化数$N$，光程差：
\begin{equation}
    \delta = 2|\Delta n| L = N\lambda
    \label{eq:fringe-count}
\end{equation}
若已知波长$\lambda$和介质长度$L$，测量条纹变化数量$N$，便可以得到折射率变化$\Delta n$：
\begin{equation}
    \Delta n = \frac{N\lambda}{2L}
    \label{eq:delta-n}
\end{equation}

因此，空气中折射率的形式：
\begin{equation}
    n = 1 + \frac{N\lambda}{2L} \cdot \frac{p}{|\Delta p|}
    \label{eq:air-n}
\end{equation}

\subsection{单缝衍射}

根据光程差公式和叠加原理，单缝衍射出现\textbf{亮条纹}的条件：
\begin{equation}
    a\sin\theta = (2k+1)\frac{\lambda}{2} \qquad (k = \pm 1, \pm 2, \ldots)
    \label{eq:bright-fringe}
\end{equation}
其中$a$为狭缝宽度，$\theta$为衍射角，$\lambda$为波长，$k$为级次。

当$\theta$很小时，近似有：
\begin{equation}
    \tan\theta = \sin\theta = \frac{x_k}{l}
    \label{eq:small-angle}
\end{equation}
其中$x_k$为第$k$级亮纹距中心亮纹的距离，$l$为缝到屏的距离。

得狭缝宽度：
\begin{equation}
    a = \frac{(2k+1) l \lambda}{2 x_k}
    \label{eq:slit-width}
\end{equation}

单缝衍射出现\textbf{暗条纹}的条件：
\begin{equation}
    a\sin\theta = k\lambda \qquad (k = \pm 1, \pm 2, \ldots)
    \label{eq:dark-fringe}
\end{equation}
相应地：
\begin{equation}
    a = \frac{k l \lambda}{x_k}
    \label{eq:slit-width-dark}
\end{equation}

\subsection{反射光栅}

狭缝光栅由大量等宽等间距的平行狭缝构成，广泛应用于各类需要分光的实验和光学仪器中。光栅方程为：
\begin{equation}
    d(\sin\theta_k \pm \sin\phi) = k\lambda
    \label{eq:grating-general}
\end{equation}
其中$d$为光栅常数，$\theta_k$为衍射角，$\phi$为入射角。对于正入射，上式简化为：
\begin{equation}
    d\sin\theta_k = k\lambda
    \label{eq:grating-normal}
\end{equation}

手机显示屏就是一个矩形网格的反射二维光栅。利用常见的激光笔作为光源，可以测量待测物品产生的衍射谱线的衍射角，再根据已知的参数，通过计算获得其他参量的数值，从而测定手机屏幕像素大小。

% ============================================================
\section{实验仪器与设备}
% ============================================================

\begin{itemize}
    \item \textbf{光源}：532\,nm半导体绿激光器、红绿蓝三色激光笔。
    \item \textbf{光学平台}：扩束镜、分光棱镜及夹具、平面镜（M1、M2）、偏振片、光阑、光学夹具、挡板、狭缝等。
    \item \textbf{气室}：长度$L = 80\,\mathrm{mm}$的玻璃气室，配气球阀门和气压表。
    \item \textbf{其他材料}：方格纸、大白板、手机自拍夹具、卷尺。
\end{itemize}

% ============================================================
\section{实验内容与步骤}
% ============================================================

\subsection{实验一：迈克尔逊干涉实验——光学平台搭建}

\begin{enumerate}
    \item \textbf{降低激光光强}：打开激光灯，调节电源旋钮调整激光强度。当光强过弱时会出现双光点，建议提升光强避免此现象。另可在激光灯前放置偏振片，略倾斜以防反射光射入激光灯，调整两偏振片的角度控制出射光强。
    \item \textbf{激光器调平行}：将光阑调整成小孔，先放在激光灯近处，调整光阑高度使激光通过小孔；然后将光阑放在激光灯远处，调整激光灯仰角使激光通过小孔。反复远近调节，使激光调平行。
    \item \textbf{平面镜调平}：
    \begin{itemize}
        \item M2镜调平：粗调旋钮使平面镜与夹具空隙基本平行。调将光阑放在M2镜近处，调整M2镜上下及左右旋钮使光通过小孔；然后将光阑放在较远处，调整C镜上下旋钮使光通过小孔。反复调节，直至M2镜调整完毕。
        \item M1镜调平：同M2镜调节方法。
    \end{itemize}
    \item \textbf{分光棱镜调平}：先粗调使空隙基本平行。渐进法细调，利用分光棱镜的玻璃反射光进行平面垂直校准。将分光棱镜放在光路上，将小孔放在较远处，调整棱镜上下旋钮使玻璃反射光通过小孔，此时棱镜一面调整完毕。反复调节两个面，使之垂直传播光路。
    \item \textbf{干涉条纹形成}：搭建光路，放置观察屏（\textbf{切勿用眼直接看}）。调整M1、M2镜的旋钮，使两反射的光斑在屏上重合。在偏振镜的后面放置扩束镜，调整扩束镜高度及位置直到干涉图案大小合适，且各方向亮度均匀，观察到类似等倾干涉的条纹。
\end{enumerate}

\subsection{实验二：观察等倾干涉并测定空气折射率}

\begin{enumerate}
    \item \textbf{等倾干涉}：微调整M1镜的上下、左右位置，使干涉圆环的圆心到达干涉图案的中心，此时M1、M2镜垂直，出现等倾干涉花样。
    \item \textbf{放置气室}：调至等倾干涉，放置空气腔，气室的长度$L=80\,\mathrm{mm}$。
    \item \textbf{充气与放气}：关闭气球的阀门，鼓气使气压值不超过气压表临界值。打开阀门，慢慢放气。记录压强变化$\Delta p$和吞吐条纹的个数$N$，利用公式~\eqref{eq:air-n}计算空气折射率。
\end{enumerate}

\subsection{实验三：观察等厚干涉}

\begin{enumerate}
    \item 调整M1镜的旋钮使同心圆环的外圈线条出现，此时M1、M2镜已不垂直。
    \item 调节M2镜距离分光镜的距离，使M1、M2镜距离分光镜距离相等，此时弯条纹变直，出现等厚干涉条纹。
\end{enumerate}

\subsection{实验四：观察衍射现象，利用激光器测定狭缝宽度}

\begin{enumerate}
    \item 采用绿色激光器，输出激光的波长为$\lambda_0 = 532\,\mathrm{nm}$。
    \item 在激光器前放置偏振片，用以调整光强。
    \item 放置狭缝，调节狭缝缝宽，使衍射效果明显。
    \item 测量缝到屏的距离$l$，在方格纸上记录可见明纹的总数目和总宽度$\Delta x$，利用暗纹间距公式计算狭缝宽度$a$。
\end{enumerate}

\subsection{实验五：利用狭缝测定激光笔的波长}

\begin{enumerate}
    \item 保持狭缝宽度和缝屏距离不变，更换待测红色激光笔。
    \item 记录红光衍射条纹的总数目和总宽度$\Delta x$。
    \item 利用实验四已测得的狭缝宽度$a$和缝屏距离$l$，由暗纹间距公式反推波长：
    \begin{equation}
        \lambda = \frac{a \cdot \Delta x}{N \cdot l}
    \end{equation}
    其中$N$为暗纹间距的数目。
\end{enumerate}

\subsection{实验六：利用反射光栅测定显示屏像素尺寸}

\begin{enumerate}
    \item 使用绿色激光（$\lambda = 532\,\mathrm{nm}$）照射显示屏（测试靶标），使其作为反射光栅。
    \item 调节入射角接近正入射，使衍射光斑投射到白屏上。
    \item 测量显示屏到白屏的距离$l$和相邻衍射光斑的间距$x$。
    \item 利用小角近似$d \approx \lambda l / x$计算像素尺寸$d$。
    \item 对不同标称分辨率（600、300、100\,PPI）的区域分别测量，比较测量值与理论值。
\end{enumerate}

% ============================================================
\section{数据处理与结果}
% ============================================================

\subsection{实验二：空气折射率测量}

\subsubsection{实验数据}

气室长度$L = 80\,\mathrm{mm}$，激光波长$\lambda = 532\,\mathrm{nm}$。

\textbf{测量一}（两点法）：记录气压从$186\,\mathrm{mmHg}$缓慢放气至$82.7\,\mathrm{mmHg}$，观察到条纹吞吐数$N = 10$个。

\textbf{测量二}（连续记录法）：每产生2个干涉圆环记录一次气压读数，得到如下数据：

\begin{table}[H]
\centering
\caption{空气折射率测量——连续记录数据}
\begin{tabular}{cccc}
\toprule
测量次序 & 气压读数 (mmHg) & 相邻气压差 (mmHg) & 累计条纹数 $N$ \\
\midrule
1 & 209 & -- & 0 \\
2 & 188 & 21 & 2 \\
3 & 166 & 22 & 4 \\
4 & 144 & 22 & 6 \\
5 & 125 & 19 & 8 \\
6 & 105 & 20 & 10 \\
7 & 83 & 22 & 12 \\
8 & 65 & 18 & 14 \\
\bottomrule
\end{tabular}
\end{table}

每产生2个条纹的气压下降平均值约为$20.6\,\mathrm{mmHg}$，数据线性度良好。

\subsubsection{折射率计算}

由式~\eqref{eq:air-n}，空气折射率：
\[
n = 1 + \frac{N\lambda}{2L} \cdot \frac{p_0}{|\Delta p|}
\]

其中$p_0 = 760\,\mathrm{mmHg}$为标准大气压（与$\Delta p$单位保持一致）。

\textbf{测量一}（$\Delta p = 186 - 82.7 = 103.3\,\mathrm{mmHg}$，$N = 10$）：
\[
\frac{N\lambda}{2L} = \frac{10 \times 532 \times 10^{-9}}{2 \times 0.08} = 3.325 \times 10^{-5}
\]
\[
n_1 = 1 + 3.325 \times 10^{-5} \times \frac{760}{103.3} = 1 + 2.446 \times 10^{-4} = 1.000245
\]

\textbf{测量二}（$\Delta p = 209 - 65 = 144\,\mathrm{mmHg}$，$N = 14$）：
\[
\frac{N\lambda}{2L} = \frac{14 \times 532 \times 10^{-9}}{2 \times 0.08} = 4.655 \times 10^{-5}
\]
\[
n_2 = 1 + 4.655 \times 10^{-5} \times \frac{760}{144} = 1 + 2.457 \times 10^{-4} = 1.000246
\]

取两次测量的平均值：
\[
\bar{n} = \frac{1.000245 + 1.000246}{2} = 1.000245
\]

文献值$n_{\text{ref}} = 1.00027821$，相对误差：
\[
E_r = \frac{|\bar{n} - n_{\text{ref}}|}{n_{\text{ref}} - 1} \times 100\% = \frac{|2.45 - 2.78| \times 10^{-4}}{2.7821 \times 10^{-4}} \times 100\% \approx 11.9\%
\]

误差主要来源于放气速度较快导致的条纹计数偏少，以及气压表读数精度有限。

\subsection{实验四：单缝衍射测狭缝宽度}

\subsubsection{实验数据}

激光波长$\lambda_0 = 532\,\mathrm{nm}$，缝屏距离$l = 760\,\mathrm{mm}$。

在光屏（方格纸）上观察到9个明纹（含中央明纹），测量第1个明纹左边缘到第9个明纹左边缘的总距离$\Delta x = 40\,\mathrm{mm}$（4\,cm）。该距离对应9个暗纹间距。

由暗纹条件~\eqref{eq:dark-fringe}，相邻暗纹间距为$\Delta x_{\text{暗}} = l\lambda/a$，故9个暗纹间距之和为：
\[
\Delta x = 9 \cdot \frac{l\lambda}{a}
\]

解得狭缝宽度：
\[
a = \frac{9 \times l \times \lambda}{\Delta x}
  = \frac{9 \times 760 \times 532 \times 10^{-6}}{40}
  = \frac{3.639}{40}
  = 0.091\,\mathrm{mm}
\]

该数值处于实验用单缝的常见规格范围（0.05\,mm $\sim$ 0.2\,mm）内，结果合理。

\subsection{实验五：测定激光笔波长}

保持狭缝宽度$a = 0.091\,\mathrm{mm}$和缝屏距离$l = 760\,\mathrm{mm}$不变，将绿色激光替换为待测红色激光笔。

在光屏上测量中央亮纹中心到单侧第4级暗纹的距离$x_4 = 21\,\mathrm{mm}$（2.1\,cm）。

由暗纹条件~\eqref{eq:dark-fringe}，$a = k\lambda l / x_k$，反解波长：
\[
\lambda_{\text{红}} = \frac{a \cdot x_4}{4 \times l}
= \frac{0.091 \times 21}{4 \times 760}
= \frac{1.911}{3040}
= 6.29 \times 10^{-4}\,\mathrm{mm}
= 629\,\mathrm{nm}
\]

红色激光笔标称波长约为$630\,\mathrm{nm}$，测量相对误差小于$0.2\%$，测量结果良好。

\subsection{实验六：反射光栅测像素尺寸}

使用绿色激光（$\lambda = 532\,\mathrm{nm}$），以显示屏作为反射光栅，屏幕至白屏距离$l = 1000\,\mathrm{mm}$。分别测量标称分辨率为600、300、100\,PPI的三个区域的衍射光斑间距。

由光栅方程~\eqref{eq:grating-normal}，取$k=1$并采用小角近似：
\[
d = \frac{\lambda \cdot l}{x}
\]

\begin{table}[H]
\centering
\caption{反射光栅测像素尺寸——测量数据与计算结果}
\begin{tabular}{ccccc}
\toprule
测试区域 (PPI) & 光斑间距 $x$ (mm) & 测量像素尺寸 $d$ ($\mu$m) & 理论像素尺寸 ($\mu$m) & 相对误差 \\
\midrule
600 & 10.0 & 53.2 & 42.3 & 25.8\% \\
300 & 5.2 & 102.3 & 84.7 & 20.8\% \\
100 & 3.0 & 177.3 & 254.0 & 30.2\% \\
\bottomrule
\end{tabular}
\end{table}

其中理论像素尺寸由$d_{\text{理论}} = 25.4\,\mathrm{mm} / \mathrm{PPI}$计算得到。

600\,PPI和300\,PPI区域的测量结果趋势正确（分辨率越高，像素尺寸越小），数据一致性较好。100\,PPI区域由于光斑过于密集（间距仅3\,mm），测量的相对误差较大。

% ============================================================
\section{分析与讨论}
% ============================================================

\subsection{等倾干涉与等厚干涉的对比}

等倾干涉条件是M1与M2平面镜严格垂直，此时等效空气薄膜为平行平面，入射角相同的光形成同一条纹，干涉图样为同心圆环。等厚干涉条件是M1与M2不垂直但距分光镜等距，等效空气薄膜为楔形，厚度相同的位置形成同一条纹，干涉图样为平行直条纹。

实验中通过微调M1镜的旋钮，可以观察到从等倾干涉（同心圆环）到等厚干涉（直条纹）的连续过渡，这直观地展示了两种干涉模式的物理联系。

\subsection{空气折射率测量的误差分析}

本实验测得空气折射率$\bar{n} = 1.000245$，与理论值$1.00027821$的相对误差约$11.9\%$。误差来源包括：

\begin{enumerate}
    \item \textbf{条纹计数偏少}：放气速度较快时，条纹吞吐速度快，人眼计数容易漏数，导致$N$偏小，进而使$n$偏小。这是主要误差来源。连续记录法（测量二）通过每2个条纹记录一次气压，有效降低了漏数风险。
    \item \textbf{气压读数误差}：手捏式气压表的最小分度为$2\,\mathrm{mmHg}$，估读精度约为$1\,\mathrm{mmHg}$。从连续记录数据看，每2个条纹的气压变化在$18\sim22\,\mathrm{mmHg}$之间波动，说明读数存在一定随机误差。
    \item \textbf{气室长度}：气室标称长度$L=80\,\mathrm{mm}$存在加工误差，但相对误差较小。
    \item \textbf{温湿度影响}：折射率与温湿度密切相关，实验中假设温湿度恒定，若实验过程中环境条件发生变化，会引入系统误差。
\end{enumerate}

\subsection{单缝衍射与反射光栅的误差来源}

\begin{enumerate}
    \item \textbf{缝屏距离测量}：用卷尺测量缝到屏的距离$l = 760\,\mathrm{mm}$，精度约为$1\,\mathrm{mm}$，相对误差约$0.13\%$。
    \item \textbf{条纹总宽度测量}：在方格纸上直接读取亮纹边缘位置，精度约为$0.5\,\mathrm{mm}$。对于总宽度$40\,\mathrm{mm}$，相对误差约$1.25\%$。高级衍射条纹亮度较弱，边缘定位困难。
    \item \textbf{小角近似}：暗纹位置公式使用了$\sin\theta \approx \tan\theta$的小角近似，对于本实验的衍射角（$< 2^{\circ}$），近似误差可忽略。
    \item \textbf{反射光栅测量}：600\,PPI和300\,PPI区域的测量误差约$20\%\sim26\%$，这可能与光斑本身有一定宽度、定位中心存在困难有关。100\,PPI区域光斑过于密集，误差最大。
\end{enumerate}

% ============================================================
\section{思考与总结}
% ============================================================

本实验通过迈克尔逊干涉仪和单缝衍射两大实验模块，系统地学习了波动光学中干涉和衍射的基本原理与测量方法。

在干涉部分，通过搭建迈克尔逊干涉仪，依次观察了等倾干涉、等厚干涉两种典型干涉图样，并利用气室改变光程差的方法测量了空气折射率。该方法的核心思想是将极微小的折射率变化转化为可数的条纹吞吐数，体现了干涉法在精密测量中的巨大优势。

在衍射部分，观察了单缝衍射的明暗条纹分布，利用已知波长的绿色激光（$532\,\mathrm{nm}$）测定了狭缝宽度（$a \approx 0.091\,\mathrm{mm}$）；然后保持狭缝不变，换用红色激光笔，反推出其波长约为$629\,\mathrm{nm}$，与标称值吻合。这种"互测"的思路展示了衍射公式的双向应用。反射光栅实验则将光栅方程应用于显示屏像素测量，验证了分辨率越高像素尺寸越小的定性规律。

实验中的主要体会是：光学实验对调节的精度要求极高。等高共轴调节、分光棱镜垂直校准、干涉条纹的优化等步骤都需要耐心细致地反复调整，任何一个环节的疏忽都可能导致无法观察到清晰的干涉或衍射图样。

\end{document}
